\subsection{抛物子代数}

命题 11. 设 \(\mathfrak{p} = \mathfrak{h} + \mathfrak{g}^{\mathrm{P}}\) 是 \(\mathfrak{g}\) 的包含 \(\mathfrak{h}\) 的子代数。下列条件等价：

\begin{enumerate}
\def\labelenumi{(\roman{enumi})}
\item
  p 包含 \((\mathfrak{g}, \mathfrak{h})\) 的一个 Borel 子代数；
\item
  存在 R 的房 C 使 \(\mathrm{P} \supset \mathrm{R}_{+}(\mathrm{C})\) ；
\item
  P 是抛物的，换言之（第六章 §1 第 7 节 定义 4），\(\mathrm{P} \cup (-\mathrm{P}) = \mathrm{R}\) 。
\end{enumerate}

条件 (i) 与 (ii) 的等价性由命题 7 得出。条件 (ii) 与 (iii) 的等价性由第六章 §1 第 7 节 命题 20 得出。

定义 2. g 的包含 h 且满足命题 11 的等价条件的子代数称为 \((\mathfrak{g}, \mathfrak{h})\) 的抛物子代数。g 的抛物子代数是指 \((\mathfrak{g}, \mathfrak{h}')\) 的抛物子代数，其中 \(h'\) 是 g 的一个分裂嘉当子代数。

这一定义立即推广到 \((\mathfrak{g}, \mathfrak{h})\) 是分裂约化李代数的情形。

注. 设 B 是 R 的一个基，b 是对应的 Borel 子代数。若 \(\Sigma \subset B\) ，用 \(Q_{\Sigma}\) 表示 \(\Sigma\) 中元素的具系数 \(\leq 0\) 的线性组合所成之根组成的集合；置 \(\mathfrak{p}(\Sigma) = \mathrm{R}_{+}(\mathrm{B}) \cup \mathrm{Q}_{\Sigma}\) 与 \(\mathfrak{p}_{\Sigma} = \mathfrak{h} \oplus \mathfrak{g}^{\mathrm{P}(\Sigma)}\) 。由第六章 §1 第 7 节 引理 3 与命题 20，\(p_{\Sigma}\) 是包含 b 的抛物子代数，且 g 的每个包含 b 的抛物子代数都是这样得到的。

引理 3. 设 V 是有限维实向量空间，S 是 \(V^{*}\) 中的一个根系，P 是 S 的抛物子集组成的集合；设 H 是诸 Ker \(\alpha\) （\(\alpha \in S\) ）组成的集合，F 是 V 关于 H 的面组成的集合（第五章 §1 第 2 节 定义 1）。

若 \(\mathrm{P} \in \mathcal{P}\) ，设 \(\overline{\mathrm{F}}(\mathrm{P})\) 是 \(v \in \mathrm{V}\) 中使 \(\alpha(v) \geq 0\) 对一切 \(\alpha \in \mathrm{P}\) 成立者组成的集合。若 \(\mathrm{F} \in \mathcal{F}\) ，设 \(\mathrm{P}(\mathrm{F})\) 是 \(\alpha \in \mathrm{R}\) 中使 \(\alpha(v) \geq 0\) 对一切 \(v \in \mathrm{F}\) 成立者组成的集合。

则 \(\mathrm{F} \mapsto \mathrm{P}(\mathrm{F})\) 是从 \(\mathcal{F}\) 到 \(\mathcal{P}\) 的一个双射；对一切 \(\mathrm{F} \in \mathcal{F}\) ，\(\overline{\mathrm{F}}(\mathrm{P}(\mathrm{F}))\) 是 \(\mathrm{F}\) 的闭包。

\begin{enumerate}
\def\labelenumi{\alph{enumi})}
\tightlist
\item
  设 \(P \in P\) 。存在 S 的一个房 C 以及一个子集 \(\Sigma\) （基 \(\mathrm{B}(\mathrm{C})\) 的），使 \(\mathrm{P} = \mathrm{S}_{+}(\mathrm{C}) \cup \mathrm{Q}\) ，其中 Q 是 \(\Sigma\) 中元素的具非正整系数的线性组合组成的集合（第六章 §1 第 7 节 命题 20）。置
\end{enumerate}

\[
\mathrm{B} (\mathrm{C}) = \{\alpha_ {1}, \dots , \alpha_ {l} \}, \Sigma = \{\alpha_ {1}, \dots , \alpha_ {m} \}.
\]

若 \(v \in \mathrm{V}\) ，我们有下列等价关系：

\[
\begin{array}{c} \alpha (v) \geq 0 \text {   for   all   } \alpha \in \mathrm{P} \\ \Longleftrightarrow \alpha_ {1} (v) \geq 0, \ldots \alpha_ {l} (v) \geq 0, \alpha_ {1} (v) \leq 0, \ldots , \alpha_ {m} (v) \leq 0 \\ \Longleftrightarrow \alpha_ {1} (v) = \dots = \alpha_ {m} (v) = 0, \alpha_ {m + 1} (v) \geq 0, \ldots , \alpha_ {l} (v) \geq 0, \end{array}
\]

故 \(\overline{\mathrm{F}}(\mathrm{P})\) 是集合

\[
\{v \in \mathrm{V} \mid \alpha_ {1} (v) = \dots = \alpha_ {m} (v) = 0, \alpha_ {m + 1} (v) > 0, \dots , \alpha_ {l} (v) > 0 \},
\]

的闭包，此集合是相对于 H 的一个面 F ，因为 S 的每个元素都是 \(\alpha_{1},\ldots,\alpha_{l}\) 的线性组合，其系数或者全 \(\geq0\) 或者全 \(\leq0\) 。此外，若 \(\beta=u_{1}\alpha_{1}+\cdots+u_{l}\alpha_{l}\in S\) ，

\[
\begin{array}{c} \beta \in \mathrm{P} (\mathrm{F}) \Longleftrightarrow u _ {m + 1} \geq 0, \ldots , u _ {l} \geq 0 \\ \Longleftrightarrow \beta \in \mathrm{S} _ {+} (\mathrm{C}) \text {   or   } (- \beta \in \mathrm{S} _ {+} (\mathrm{C}) \text {   and   } u _ {m + 1} = \ldots = u _ {l} = 0) \\ \Longleftrightarrow \beta \in \mathrm{S} _ {+} (\mathrm{C}) \cup \mathrm{Q} = \mathrm{P}, \end{array}
\]

故 P(F) = P。

\begin{enumerate}
\def\labelenumi{\alph{enumi})}
\setcounter{enumi}{1}
\tightlist
\item
  设 \(F \in F\) 。显然 \(P(F) \in \mathcal{P}\) 。另一方面，F 含于相对于 \(\mathcal{H}(\text{Chap. V}, \S1, \text{no. 3, formulas (6)}),\) 的一个房的闭包之中，因而是相对于该房的壁所成集合的一个面（第五章 \(\S1\) 第 4 节 命题 9）。于是，\(\overline{F}\) 具有形状 \(\{v \in V \mid \alpha(v) \geq 0 \text{ 对一切 } \alpha \in T\}\) ，其中 T 是 S 的一个子集，而显然可取它等于 \(P(F)\) 。这样，\(\overline{F} = \overline{F}(P(F))\) 。证毕。
\end{enumerate}

若 \(P \in P\) ，使 \(P = P(F)\) 的面 F 称为与 P 相伴的面；我们把它记作 \(F(P)\) 。我们把这些约定推广到 \((\mathfrak{g}, \mathfrak{h})\) 是分裂约化的情形。

命题 12. 设 \(\mathcal{H}\) 是 \(\mathfrak{h}_{\mathbf{R}}\) 的超平面组成的集合，这些超平面由 R 中根的核构成。设 \(\mathcal{F}\) 是 \(\mathfrak{h}_{\mathbf{R}}\) 相对于 \(\mathcal{H}\) 的面组成的集合。设 \(\mathcal{S}\) 是 \((\mathfrak{g},\mathfrak{h})\) 的抛物子代数组成的集合。对每个 \(\mathfrak{p} = \mathfrak{h} + \mathfrak{g}^{\mathrm{P}}\in \mathcal{S}\) ，设 \(\mathrm{F}(\mathfrak{p})\) 是与 P 相伴的面。则 \(\mathfrak{p}\mapsto \mathrm{F}(\mathfrak{p})\) 是从 \(\mathcal{S}\) 到 \(\mathcal{F}\) 的一个双射。若 \(\mathfrak{p}_1,\mathfrak{p}_2\in \mathcal{P}\) ，

\[
\mathfrak {p} _ {1} \supset \mathfrak {p} _ {2} \Longleftrightarrow \mathrm{F} (\mathfrak {p} _ {1}) \subset \overline {{\mathrm{F} (\mathfrak {p} _ {2})}}.
\]

这由引理 3 立即得出。

例. 对应于 \((\mathfrak{g}, \mathfrak{h})\) 的包含一个 Borel 代数 b 的抛物子代数的那些面，就是含于与 b 相伴的房的闭包之中的那些面（参见上面的注）。

命题 13. 设 \(\mathfrak{p} = \mathfrak{h} + \mathfrak{g}^{\mathrm{P}}\) 是 \((\mathfrak{g},\mathfrak{h})\) 的一个抛物子代数，Q 是 \(\alpha \in \mathrm{P}\) 中使 \(-\alpha \notin \mathrm{P}\) 者组成的集合，并设 \(\mathfrak{s} = \mathfrak{h} + \mathfrak{g}^{\mathrm{P}\cap (-\mathrm{P})}\) 。则 \(\mathfrak{p} = \mathfrak{s} \oplus \mathfrak{g}^{\mathrm{Q}}\) ，\(\mathfrak{s}\) 在 \(\mathfrak{g}\) 中是约化的，且 \(\mathfrak{g}^{\mathrm{Q}}\) 是 \(\mathfrak{p}\) 的最大幂零理想并且是 \(\mathfrak{p}\) 的幂零根基。\(\mathfrak{p}\) 的中心为零。

由命题 2，\(\mathfrak{s}\) 在 \(\mathfrak{g}\) 中是约化的，且 \(\mathfrak{g}^{\mathrm{Q}}\) 是 \(\mathfrak{p}\) 的一个幂零理想。若 \(\mathfrak{n}\) 是 \(\mathfrak{p}\) 的最大幂零理想，则 \(\mathfrak{g}^{\mathrm{Q}} \subset \mathfrak{n} \subset \mathfrak{h} + \mathfrak{g}^{\mathrm{Q}}\) （命题 2 (i)）；若 \(x \in \mathfrak{n} \cap \mathfrak{h}\) ，则 \(\operatorname{ad}_{\mathfrak{p}} x\) 是幂零的，故 \(\alpha(x) = 0\) 对一切 \(\alpha \in \mathrm{P}\) 成立，从而 \(x = 0\) ；这就证明了 \(\mathfrak{n} = \mathfrak{g}^{\mathrm{Q}}\) 。由于 \([\mathfrak{h}, \mathfrak{g}^{\mathrm{Q}}] = \mathfrak{g}^{\mathrm{Q}}\) ，\(\mathfrak{p}\) 的幂零根基包含 \(\mathfrak{g}^{\mathrm{Q}}\) ，从而等于 \(\mathfrak{g}^{\mathrm{Q}}\) 。设 \(z = h + \sum_{\alpha \in \mathrm{P}} u_{\alpha}\) （其中 \(h \in \mathfrak{h}, u_{\alpha} \in \mathfrak{g}^{\alpha}\) ）是 \(\mathfrak{p}\) 的中心的一个元素。对一切 \(h' \in \mathfrak{h}\) ，有 \(0 = [h', z] = \sum \alpha(h')u_{\alpha}\) ，故 \(u_{\alpha} = 0\) 对一切 \(\alpha \in \mathrm{P}\) 成立；由此推出 \([h, \mathfrak{g}^{\beta}] = 0\) 对一切 \(\beta \in \mathrm{P}\) 成立，故 \(h = 0\) 。

